
\section{Proof of the general Gram normal form}
\label{app:normal-form-proof}

This appendix proves Lemma~\ref{lem:normal-form}.  We follow the three
statements of the lemma in order.

\begin{proof}[Proof of Lemma~\ref{lem:normal-form}]
Put $r=\binom m2$.

\medskip
\noindent
\textbf{1. Reduction of the Gram family.}
By Lemma~\ref{lem:frame}, every Gram representation of $\calF_m$
differs from the fixed representation
\[
G_m=
\begin{pmatrix}
D&K&0\\
K&D&0\\
0&0&B
\end{pmatrix}
\]
by Pl\"ucker corrections.  We classify these corrections according to
how their four labels are distributed between the two sides of the
Toeplitz corner.

Four labels on the same side produce corrections acting entirely within
one pure wedge family.  Such a correction has entries
$
(\mathcal E_0)_{ab,cd}=\Omega_{abcd},
$
where $\Omega$ is totally antisymmetric in its four indices.  These are
the same-side Pl\"ucker corrections, or equivalently the pure
four-forms.

Two labels on each side produce the Pl\"ucker relations that couple the
two pure families with the mixed family.  After the pure cross
correction is collected into a symmetric matrix
$
\mathcal E_1\in\mathbb S^r,
$
the corresponding change in the mixed block is determined by the same
Pl\"ucker relations.  We denote this induced matrix by
$
\calR\mathcal E_1\in\mathbb S^{m^2}.
$
Equivalently, after antisymmetrically extending the entries of
$\mathcal E_1$ in each pure index pair, the realignment map is given by
\[
 (\calR\mathcal E_1)_{ac,bd}
 =(\mathcal E_1)_{ab,cd},
 \qquad
 (\calR\mathcal E_1)_{ad,bc}
 =-(\mathcal E_1)_{ab,cd}.
\]
Thus $\calR$ simply re-pairs the four indices from the two pure wedge
pairs into two mixed wedge pairs.

Corrections involving three labels on one side and one label on the
other produce pure--mixed blocks.  To remove these terms, change the
sign of all variables on the negative side.  This transformation leaves
the pure coordinates $u,v$ unchanged and changes the sign of every
mixed coordinate $w$.  The polynomial $\calF_m$ is unchanged, so if
$Q\succeq0$ is a Gram representation, averaging $Q$ with its image
under this symmetry preserves both the represented polynomial and
positive semidefiniteness, while eliminating all pure--mixed blocks.

Next interchange the two sides of the Toeplitz corner.  Under this
symmetry,
\[
                         (u,v)\longmapsto(-v,-u),
 \qquad
                         w_{ab}\longmapsto-w_{ba}.
\]
A second averaging therefore makes the two pure diagonal corrections
equal and makes the pure cross correction symmetric.  Consequently,
whenever a positive semidefinite Gram representation exists, there is
one of the form
\[
 Q_{\rm av}
 =
 G_m+
 \begin{pmatrix}
 \mathcal E_0&\mathcal E_1&0\\
 \mathcal E_1&\mathcal E_0&0\\
 0&0&\calR\mathcal E_1
 \end{pmatrix},
\]
which is exactly \eqref{eq:general-gram-1}.  Both averaging operations
are convex averages of positive semidefinite matrices and therefore
preserve positive semidefiniteness.

\medskip
\noindent
\textbf{2. Block diagonalization.}
Write
\[
 A:=D+\mathcal E_0,
 \qquad
 C:=K+\mathcal E_1.
\]
Then the two pure blocks of $Q_{\rm av}$ have the form
\[
 \begin{pmatrix}
 A&C\\
 C&A
 \end{pmatrix}.
\]
Consider the orthogonal matrix
\[
 U=
 \begin{pmatrix}
 \frac1{\sqrt2}I_r&\frac1{\sqrt2}I_r&0\\
 \frac1{\sqrt2}I_r&-\frac1{\sqrt2}I_r&0\\
 0&0&I_{m^2}
 \end{pmatrix}.
\]
It corresponds to the change of coordinates
\[
 (u,v,w)
 \longmapsto
 \left(
 \frac{u+v}{\sqrt2},
 \frac{u-v}{\sqrt2},
 w
 \right).
\]
Since $U^\top U=I$, direct multiplication gives
\[
\begin{aligned}
 U^\top Q_{\rm av}U
 &=
 \begin{pmatrix}
 A+C&0&0\\
 0&A-C&0\\
 0&0&B+\calR\mathcal E_1
 \end{pmatrix}                                                   \\[1mm]
 &=
 \begin{pmatrix}
 D+\mathcal E_0+K+\mathcal E_1&0&0\\
 0&D+\mathcal E_0-K-\mathcal E_1&0\\
 0&0&B+\calR\mathcal E_1
 \end{pmatrix}.
\end{aligned}
\]
Hence
\[
 P_+=D+\mathcal E_0+K+\mathcal E_1,\qquad
 P_-=D+\mathcal E_0-K-\mathcal E_1,\qquad
 M=B+\calR\mathcal E_1,
\]
which proves \eqref{eq:general-gram-2}--\eqref{eq:sectors}.

\medskip
\noindent
\textbf{3. Equivalence with the SoS condition.}
Suppose first that $\calF_m$ is SoS.  By Lemma~\ref{lem:frame}, it has a
positive semidefinite Gram representation in the restricted wedge
coordinates.  Part~1 shows that this representation can be replaced,
without losing positive semidefiniteness, by one of the form
\eqref{eq:general-gram-1}.  Thus there exist
$\mathcal E_0,\mathcal E_1$ for which $Q_{\rm av}\succeq0$.

Conversely, if such matrices $\mathcal E_0,\mathcal E_1$ exist and
$Q_{\rm av}\succeq0$, then
$
Q_{\rm av}=R^\top R
$
for some real matrix $R$.  Since $Q_{\rm av}$ is a Gram representation
of $\calF_m$,
\[
 \calF_m
 =z_r^\top Q_{\rm av}z_r
 =\|Rz_r\|^2,
\]
which is a sum of squares of homogeneous quadratic polynomials.
Therefore
\[
 \calF_m\text{ is SoS}
 \quad\Longleftrightarrow\quad
 \exists\,\mathcal E_0,\mathcal E_1
 \text{ such that }Q_{\rm av}\succeq0.
\]

Finally, $U$ is orthogonal, so
\[
 Q_{\rm av}\succeq0
 \quad\Longleftrightarrow\quad
 U^\top Q_{\rm av}U\succeq0.
\]
By the block diagonal form in Part~2, this is equivalent to
\[
                         P_+\succeq0,\qquad
                         P_-\succeq0,\qquad
                         M\succeq0.
\]
This proves the lemma.
\end{proof}